\section*{Chapter \ref{ch:b3:np-completeness} --- NP-Completeness}

\begin{solution}{exo:b3:np-completeness:1}
$S \le_p \mathrm{SAT}$ proves that $S$ is in NP, no harder than SAT; nothing about its hardness. For NP-hardness the reduction must go the other way: from SAT, or another known \omterm{def:b3:np-completeness:reduction}{NP-complete} problem, to $S$. With both, $S$ is \omterm{def:b3:np-completeness:reduction}{NP-complete}.
\end{solution}

\begin{solution}{exo:b3:np-completeness:2}
$(x_1 \vee \neg x_2 \vee y_1)(\neg y_1 \vee x_3 \vee y_2)(\neg y_2 \vee x_4 \vee \neg x_5)$. With $x_4$ true and the others false, the first clause needs $y_1$ true, the second then needs $y_2$ true, and the third is satisfied by $x_4$ (and by $\neg x_5$). The rule: the $y$s before the clause of the true literal are true, those after are false.
\end{solution}

\begin{solution}{exo:b3:np-completeness:3}
Six vertices: a triangle $x_1, x_2, \neg x_3$ and a triangle $\neg x_1, x_3, x_2$, plus the edges $x_1 - \neg x_1$ and $\neg x_3 - x_3$. The two $x_2$ vertices are not joined, so $\{x_2, x_2'\}$, one per triangle, is independent of size 2: the assignment $x_2$ true, the others arbitrary. Another is $\{x_1, x_3\}$: $x_1$ and $x_3$ true.
\end{solution}

\begin{solution}{exo:b3:np-completeness:4}
Digits $x_1 x_2 x_3 \mid c_1$: $y_1 = 1001$, $z_1 = 1000$, $y_2 = 100$, $z_2 = 101$, $y_3 = 11$, $z_3 = 10$, two slacks $1$ and $1$; target $1113$. For example $y_1 + y_2 + y_3 + 1 = 1001 + 100 + 11 + 1 = 1113$: all three variables true; the clause digit gets 1 from $y_1$, 1 from $y_3$ and 1 from a slack.
\end{solution}

\begin{solution}{exo:b3:np-completeness:5}
Joining literals of the same clause lets a \omterm{def:b3:np-completeness:problems}{clique} take two literals from one clause and none from another. The eight clauses that forbid every assignment of $x_1, x_2, x_3$ form an unsatisfiable formula, yet the buggy graph has a \omterm{def:b3:np-completeness:problems}{clique} of size 8: the literal $x_1$ occurs in four clauses and $x_2$ in four, and these eight occurrences are pairwise non-complementary, so all joined; a test checks it. The fix joins two occurrences only if they lie in different clauses and are not complementary; then a \omterm{def:b3:np-completeness:problems}{clique} of size $m$ takes one literal per clause, consistently.
\end{solution}

\begin{solution}{exo:b3:np-completeness:6}
\code{exercises/ex6_col_to_sat.py}: variables $x_{v,c}$ for 3 colours; per vertex one clause $x_{v,0} \vee x_{v,1} \vee x_{v,2}$ and three clauses $\neg x_{v,a} \vee \neg x_{v,b}$; per edge $uv$ and colour $c$, $\neg x_{u,c} \vee \neg x_{v,c}$. A colouring gives a satisfying assignment and back; $3n$ variables and $4n + 3m$ clauses (\code{test_exercise_6_colouring_to_sat}).
\end{solution}

\begin{solution}{exo:b3:np-completeness:7}
Vertex cover is in NP: the set is a certificate, checked edge by edge. Facts used: $S$ is independent if and only if $V \setminus S$ covers every edge, since an edge with no end in $V \setminus S$ has both ends in $S$. Map $(G, k)$ to $(G, n - k)$: $G$ has an \omterm{def:b3:np-completeness:problems}{independent set} of size $k$ if and only if it has a vertex cover of size $n - k$. The map is computed in linear time, so \omterm{def:b3:np-completeness:problems}{independent set} $\le_p$ vertex cover, and vertex cover is \omterm{def:b3:np-completeness:reduction}{NP-complete}.
\end{solution}

\begin{solution}{exo:b3:np-completeness:8}
Cell variables: $7(t + 1)W$. Window variables: one per legal window and per position, $t(W - 2)L$. Clauses: per cell, one ``at least one'' and $\binom{7}{2} = 21$ ``at most one'': $22(t + 1)W$; the start row, $W$; acceptance, 1; per window position, $6L + 1$: $t(W - 2)(6L + 1)$. For $t = 3$, $W = 7$, $L = 193$: $196 + 15 \times 193 = 3{,}091$ variables and $616 + 7 + 1 + 15 \times 1{,}159 = 18{,}009$ clauses, as printed (\code{test_cook_levin_counts_and_windows}).
\end{solution}

\begin{solution}{exo:b3:np-completeness:9}
Facts used: the graph is symmetric, $a \to b$ being an edge if and only if $\neg b \to \neg a$ is, so a path from $a$ to $b$ gives one from $\neg b$ to $\neg a$, and the components come in mirror pairs; a satisfying assignment never has a true literal implying a false one along an edge.

If $x$ and $\neg x$ share a component, $x \Rightarrow \neg x$ and $\neg x \Rightarrow x$ along paths, so neither value is possible. Conversely, number the components in a topological order and make $x$ true when its component comes after that of $\neg x$. Suppose a clause $a \vee b$ is false: $a$ and $b$ are false, so the component of $\neg a$ comes after that of $a$, and the same for $b$. The edge $\neg a \to b$ puts $b$'s component at or after $\neg a$'s, and the edge $\neg b \to a$ puts $a$'s at or after $\neg b$'s. Then $a \le \neg a \le b \le \neg b \le a$ in the order, with $a$ strictly before $\neg a$: a contradiction. The algorithm finds the components and their order in linear time (\code{test_two_sat}).
\end{solution}

\begin{solution}{exo:b3:np-completeness:10}
Facts used: in a 3-colouring the palette T, F, B takes three different colours, which we name by it; a literal vertex is adjacent to B, so it is T or F, and $x$, $\neg x$ differ.

The OR \omterm{def:b3:np-completeness:gadget}{gadget} on inputs $a$, $b$ coloured T or F: if both are F, $u$ and $v$ avoid F, so they are T and B in some order, and the output, adjacent to both, is F. If $a$ is T, colour $u$ F, $v$ B (allowed, since $b$ is T or F) and the output T; if $b$ is T, symmetrically. So the output can be T exactly when $a$ or $b$ is T, and when it is not T it is F.

If the graph is 3-colourable, read the assignment from the literal vertices. Each clause's last output is adjacent to F and B, so it is T; then the first output or $c$ is T, and then $a$, $b$ or $c$ is: the clause is satisfied. Conversely, from a satisfying assignment colour the literals, then each \omterm{def:b3:np-completeness:gadget}{gadget} as above with output T whenever possible; the last output of each clause is T, compatible with its edges to F and B. Sizes are linear, so the reduction is polynomial.
\end{solution}

\begin{solution}{exo:b3:np-completeness:11}
Subset sum is in NP: the subset is the certificate. The reduction of the chapter is correct. If the formula is satisfiable, take $y_i$ or $z_i$ as $x_i$ is true or false; each variable digit is 1, each clause digit has 1 to 3 from the literals, completed to 3 by slacks. If a subset reaches the target, no digit carries, since a digit's total is at most 5; each variable digit is 1, so exactly one of $y_i$, $z_i$ is chosen, an assignment; each clause digit is 3 with at most 2 from slacks, so some true literal is in the clause. The numbers have $n + m$ digits, written in polynomial time. The dynamic program takes $O(nT)$ steps, and here $T \approx 10^{n + m}$: exponential in the input size, which is about $(2n + 2m)(n + m)$ digits.
\end{solution}

\begin{solution}{exo:b3:np-completeness:12}
Let $C$ be \omterm{def:b3:np-completeness:reduction}{NP-complete} and in \omterm{def:b3:np-completeness:conp}{co-NP}. For $A$ in NP, $A \le_p C$ by a function $f$; then $\overline{A} \le_p \overline{C}$ by the same $f$, and $\overline{C}$ is in NP, so $\overline{A}$ is in NP: $A$ is in \omterm{def:b3:np-completeness:conp}{co-NP}. So $\mathrm{NP} \subseteq \mathrm{co\text{-}NP}$. For $B$ in \omterm{def:b3:np-completeness:conp}{co-NP}, $\overline{B}$ is in NP, hence in \omterm{def:b3:np-completeness:conp}{co-NP}, so $B$ is in NP: $\mathrm{co\text{-}NP} \subseteq \mathrm{NP}$.
\end{solution}

\begin{solution}{pb:b3:np-completeness:1}
\textbf{1.} Given a graph and $s$, is there a map from vertices to $\{1, \ldots, s\}$ giving different values to the two ends of every edge?

\textbf{2.} The certificate is the timetable; checking looks at every vertex and edge once, $O(n + m)$.

\textbf{3.} Three: ABC and DEF are triangles; A, B, C get three slots, D a slot other than C's, E and F the other two.

\textbf{4.} Its three exams pairwise clash, so they need three different slots.

\textbf{5.} An algorithm for the fewest slots answers the decision with $s = 3$ by comparing; so the optimisation is \omterm{def:b3:np-completeness:reduction}{NP-hard} too.

\textbf{6.} The three palette vertices are pairwise adjacent, so they take three different colours; calling those colours T, F and B is just naming them.

\textbf{7.} $x$ and $\neg x$ are adjacent to B and to each other, so they take the two colours other than B's, one each.

\textbf{8.} $u$ is adjacent to $a$ (F) and $v$ to $b$ (F), so $u, v \in \{T, B\}$, different from each other; the output is adjacent to both, so it is F.

\textbf{9.} Colour $u$ F (its neighbour $a$ is T), $v$ B ($b$ is T or F), the output T.

\textbf{10.} It is adjacent to F and B.

\textbf{11.} The first \omterm{def:b3:np-completeness:gadget}{gadget}'s output is F (question 8), and then the second \omterm{def:b3:np-completeness:gadget}{gadget}'s inputs are both F, so its output is F, but it must be T.

\textbf{12.} 3 palette vertices, $2n$ literal vertices, 6 per clause.

\textbf{13.} 3 palette edges, 3 per variable, 12 per clause: 5 per OR \omterm{def:b3:np-completeness:gadget}{gadget} and 2 to F and B.

\textbf{14.} $n = 3$, $m = 2$: $6 + 12 + 3 = 21$ and $9 + 24 + 3 = 36$, as printed.

\textbf{15.} $200 + 2{,}580 + 3 = 2{,}783$ exams and $300 + 5{,}160 + 3 = 5{,}463$ clashes.

\textbf{16.} It writes $O(n + m)$ vertices and edges, each by a fixed rule.

\textbf{17.} They are linear: a formula of size $N$ becomes a graph of size $O(N)$, so if 3-SAT needs time $2^{\Omega(N)}$ in the worst case, so does 3-colouring, up to the constant in the exponent.

\textbf{18.} In a 2-colouring colours alternate along any walk, so a cycle has even length. Conversely, if there is no odd cycle, colour each component by the parity of the distance from a start vertex; an edge between two vertices of the same parity would close an odd cycle.

\textbf{19.} Breadth-first search, colouring each new vertex opposite to its parent; on an edge whose ends have the same colour, stop: the two paths to the start and that edge form an odd closed walk, which contains an odd cycle, the certificate that no timetable exists.

\textbf{20.} $x_v$ means ``exam $v$ in slot 1''; each clash $uv$ gives $(x_u \vee x_v)$ and $(\neg x_u \vee \neg x_v)$: exactly one of the two in slot 1.

\textbf{21.} With three slots a vertex needs more than one Boolean variable, and ``different slots'' is no longer a conjunction of two-literal clauses: ``at least one slot'' needs three literals.

\textbf{22.} No. It means no known algorithm is fast on every graph; with 40 exams, backtracking or a SAT solver (\cref{ch:b3:hard-problems}) finds a timetable or proves there is none quickly in practice.

\textbf{23.} That a solver found no timetable, which is not a short certificate: proofs of non-colourability can be long, since 3-colouring's complement is in \omterm{def:b3:np-completeness:conp}{co-NP} and not known to be in NP. A small certificate exists in special cases, such as four pairwise clashing exams.

\textbf{24.} Two slots: linear time, by breadth-first search or \omterm{def:b3:np-completeness:2sat}{2-SAT}, with an odd cycle as a short proof of impossibility. Three slots: \omterm{def:b3:np-completeness:reduction}{NP-complete} by the \omterm{def:b3:np-completeness:gadget}{gadgets}, with $2n + 6m + 3$ vertices and $3n + 12m + 3$ edges per formula; solved in practice by search, with no short proof of impossibility known in general.
\end{solution}
